\newpage
\addcontentsline{toc}{chapter}{\quad \: V. \hspace{9pt} CONCLUSION}

\begin{center}
  \vspace*{1.65in}
  CHAPTER V\\
  CONCLUSION   
\end{center}

\bigskip
We investigated a low-memory, FPGA-targeted VHDL architecture for $M$-QAM-SRRC SR--CF detection.  We began by defining the streaming SCA's scope and then described its sample-by-sample datapath.  The datapath does not instantiate an FFT or vendor-specific FFT intellectual-property core, although MATLAB and the Vivado toolchain were used for modeling, simulation, and synthesis.  Finally, we presented thirteen single-trial VHDL simulations and the archived resource counts.  Recovered 2021 component design reports and project scripts identify the intended development platform as a ZedBoard XC7Z020CLG484-1 with Vivado 2019.2.1.  However, the exact thesis-era RTL revision and top-level synthesis reports are unavailable, later target metadata conflicts, and the recovered 2022 checkout contains the fixed-point and flow-control discrepancies documented in Chapter III.  The archived simulations therefore cannot establish bit-true conformance, and the experiment-level resource counts cannot be independently tied to the documented component setup.  None of these results are measurements from deployed flight hardware.

% preprocessing
As discussed in Section 3.2, one shortcoming of this design is that its serial processing time grows as the SR and CF candidate grids become finer.  Koch first forms spectral segments from a thresholded Welch PSD \cite{koch-19}.  After centering and dynamically low-pass filtering each segment, the cyclic-autocorrelation FFT (CAC-FFT) forms $z[n]=|x[n]|^2$ and uses an $N$-point FFT to evaluate the zero-lag CAC on a uniformly spaced SR grid; peak finding selects candidate rates.  For each selected rate, Koch's implemented carrier-frequency detector then sweeps that segment across frequency offsets, applies the rate-dependent low-pass filter, evaluates a single-rate CAC, and thresholds the normalized peak.  A separate center-frequency preprocessor in the report was preliminary and was not used.  The hardware described in this thesis serializes and simplifies the carrier sweep, but it integrates neither the preceding PSD segmentation nor the CAC-FFT candidate-selection stage.  Adding those stages is a possible extension; it would require additional filtering, block storage, and an FFT, and the reduction in serial work would depend on how many candidates survive preprocessing, so no fixed reduction factor is established here.

% decimation
Decimation is another possible way to reduce the sample rate processed by a candidate branch, but the factor is not generally $F_s/R_s$.  After the band of interest is centered, an anti-alias low-pass filter must precede integer downsampling by $D$, giving $F_s'=F_s/D$.  The value of $D$ must preserve the occupied bandwidth and the cyclic frequencies needed by the detector under the new Nyquist limit.  It also changes the NCO control-word mapping, the observation interval represented by $N$ samples, and the scheduling of parallel streaming-SCA instances.  This thesis does not implement or evaluate such a decimator.

% postprocessing
The constant $\beta=10^{-4}$ used in the characterization was selected empirically; the thirteen single-trial plots do not calibrate a false-alarm probability or establish a probability of detection.  A recovered exploratory ROC notebook supplies no additional evidence because its stored run terminates with an execution error rather than a completed sweep.  The present row-maximum rule can report at most one CF for each SR candidate, so it cannot resolve two signals that share an SR but have different CFs.  Detecting multiple local CF maxima above threshold, followed by suitable nonmaximum suppression or clustering, is one possible extension.  Similarly, a fine SR grid can yield a cluster of correlated above-threshold candidates around one physical feature.  Those adjacent grid responses are not independent statistical false alarms and should be consolidated before a waveform count is reported.  Selecting and validating such postprocessing requires Monte Carlo receiver-operating-characteristic analysis over the intended channels and signal classes.

% additional waveforms
Lastly, the characterization tests only 4-QAM-SRRC and 16-QAM-SRRC signals.  The same datapath can evaluate its statistic for other sampled waveforms, but that fact alone does not establish useful cyclic features or reliable detection.  Rectangular pulse shapes, FSK, OOK, MSK, DSSS, DPSK, and Manchester coding were not tested here; their candidate grids, filtering, and thresholds would have to be derived and validated before support could be claimed.

\clearpage

